% Methods — SIREN
\section{Methods}

\subsection{Problem setting}

Let $X \in \mathbb{R}^{n \times p}$ be a descriptor matrix and
$y \in \mathbb{R}^{n}$ (regression) or $y \in \{0,1\}^{n}$ (classification) the
endpoint. A selector is a map
$\mathcal{S}(X, y, k) \mapsto \{j_1, \dots, j_k\}$ that returns $k$ column
indices. We consider only selectors that are \emph{refitted} on the development
rows of every evaluation split: for a split with training rows $D$ and held-out
rows $T$, the selection is $\mathcal{S}(X_D, y_D, k)$ and the score is computed
on $T$ using only the selected columns. No row of $T$ influences the selection.

\subsection{The marginal scorecard}

The reference procedure, which we call the scorecard, ranks every descriptor by
its absolute Spearman correlation with the endpoint,
$s_j = |\rho_s(X_{\cdot j}, y)|$, keeps the top candidates, and then greedily
removes any descriptor whose absolute Pearson correlation with an already kept
descriptor exceeds a threshold $\tau$ (0.8 in our runs). It has four structural
properties that motivate this work: it cannot promote a descriptor whose
univariate association is masked by a correlated partner (a \emph{suppressor});
it cannot promote an interaction partner that carries no marginal signal; it
does not account for the resampling variability of the ranking; and the number
of retained descriptors is controlled only by $\tau$ and by the user's budget.

\subsection{The SIREN objective}

SIREN replaces the marginal score with a greedy forward selection that scores
each remaining candidate $j$ \emph{conditional} on the set $S$ already selected:
\begin{equation}
J(j \mid S) = \alpha\, R(j \mid S) + \beta\, Y(j \mid S)
            + \gamma\, T(j) - \delta\, D(j \mid S),
\end{equation}
where $R$ is a conditional relevance term, $Y$ an interaction term, $T$ a
stability term and $D$ a redundancy penalty. The default weights are
$\alpha = 1$, $\beta = 0.5$, $\gamma = 0.5$ and $\delta = 0$; the redundancy
penalty is disabled because the near-duplicate screen below already enforces the
hard constraint, and a soft penalty was found to trade suppressor retention for
duplicates that the screen had already removed.

\paragraph{Conditional relevance.}
For regression, both the candidate and the endpoint are residualised against the
selected columns, $r_j = X_{\cdot j} - P_S X_{\cdot j}$ and
$r_y = y - P_S y$, and $R(j \mid S) = \hat{I}(r_j; r_y)$. For classification the
endpoint is not continuous, so only the descriptor is residualised and the label
is kept discrete, $R(j \mid S) = \hat{I}(r_j; y)$. This distinction matters: a
continuous residual of a binary label silently switches the estimator to a
regression variant whose output is not bounded by the label entropy $H(y)$, and
the resulting value cannot be compared with the marginal term or with the
nats-scale interaction term. With the label kept discrete, $R \in [0, H(y)]$ and
the suppression gain $R(j\mid S) - \hat{I}(X_{\cdot j}; y)$ has a direct reading
in bits or nats.

\paragraph{Interaction term.}
For a candidate $j$ and each selected partner $l$, we form the standardised
product $u = \tilde{X}_{\cdot j} \odot \tilde{X}_{\cdot l}$, residualise it
against the two raw columns, and measure
$Y(j \mid S) = \max_{l \in S} \hat{I}(\tilde{u}; y)$. Residualising the product
removes the squared marginal component that makes a duplicate look synergistic.
Near-duplicate pairs ($|\rho| \ge 0.999$) have no beyond-marginal signal and are
assigned $Y = 0$. The term is the absolute interaction information in nats, on the
same scale as $R$. An earlier normalised variant, which divided the lift by the
pair's marginal information, is unbounded when both marginals are near zero ---
the interaction-only regime the term exists to detect --- and on real descriptor
matrices it produced objective values in the thousands while selecting noise.

\paragraph{Stability.}
$T(j)$ is the fraction of random subsamples (80\% of rows, 16 draws by default)
in which the sign of $\hat{\rho}_s(X_{\cdot j}, y)$ matches the full-sample sign.
A descriptor whose association flips sign under resampling receives no bonus.

\paragraph{Redundancy.}
$D(j \mid S) = \max\big(0, \hat{I}(X_{\cdot j}; y) - R(j \mid S)\big)
 + \max_{l \in S} |\rho_{\text{partial}}(X_{\cdot j}, X_{\cdot l} \mid S)|$,
both terms bounded in $[0, 1]$. It is reported but not applied by default.

\subsection{Screening}

Two constraints are applied before selection and are not part of the objective.
A \emph{support} mask removes descriptors that are constant or almost constant
on the development rows (default: fewer than five distinct values). A
\emph{near-duplicate} mask removes a column whose absolute Pearson correlation
with an already available column exceeds $0.999$; the first occurrence in column
order is kept, so the result is deterministic. Because the mask is applied before
the greedy loop, a duplicated descriptor cannot consume one of the $k$ slots.

\subsection{Synergy seeding}

Greedy forward selection can only reward an interaction against a partner that
has \emph{already} been selected, so a design whose only signal is a pure
interaction is unreachable: the first step has no partner and the marginal
scores are at chance. SIREN therefore checks, before the loop, whether any
eligible candidate carries marginal information. The check is calibrated by
permutation: the endpoint is shuffled and the maximum marginal estimate over all
eligible columns defines a noise floor, and the screen fires only when
$\max_j \hat{I}(X_{\cdot j}; y) < \max(0.15, 1.5 \times \text{floor})$.
When it fires, candidate pairs are scored by their interaction information and
the strongest pair is admitted first, subject to two guards: the interaction must
reach 0.15 nats, and it must exceed the median of the screened pairs by three
median absolute deviations. On the molecular benchmarks the weakest endpoint
already carries 0.21 nats of marginal information, so the screen is skipped and
the additional cost is not paid.

\subsection{Estimators and reproducibility}

All mutual-information values are estimated with the $k$-nearest-neighbour
estimators of \textsc{scikit-learn}, with a fixed random seed. Continuous
variables are quantile-binned into 12 bins before estimation; the label is left
discrete for classification. Spearman correlations are computed with mid-ranks.
Every random operation (stability subsamples, permutation floor, pair screening)
derives its generator from the user-supplied seed, so a run is reproducible. The
selector is $O(k \cdot p)$ batched estimator calls: all candidates are scored
against the current context in one linear solve and one or two estimator calls,
which is what makes the nested protocol of the next section affordable.

\subsection{Evaluation protocol}

All selectors are compared under identical splits. For the molecular benchmarks
we use the scaffold split of each dataset (80/10/10), repeated over ten seeds; the
descriptor set is the 210 RDKit 2D descriptors after the release-time contract.
The headline metric is the held-out Spearman correlation of an ordinary
least-squares readout fitted on the selected columns, deliberately simple so that
differences reflect selection rather than model capacity. For classification the
readout is logistic regression and the metric is AUROC. Differences against the
scorecard are reported as paired bootstrap intervals over identical splits, and
family-wise error is controlled with Holm--Bonferroni within each table. A
separate experiment replaces the linear readout with a random forest on
scaffold-held-out rows to test whether the selected subspace helps a nonlinear
model.
